\Article{11}
Man denke sich von jedem ganzen Punkte $k$ nach $k+\frac12i$ gerade Linien gezogen, deren rechte Seite als positiv, die linke als negativ angesehen wird. Es sei $l$ eine Linie, und $Sl$ bezeichne die Summe aller Schnitte der $l$ mit jenem System von Linien, diejenigen als positiv angesehen, wo $l$ von der negativen auf die positive übergeht, die entgegengesetzten Schnitte als negativ. Man hat dann für den Decidenten folgenden Ausdruck
\[
\Sigma(\Upsilon l\pdot Sl)+\Sigma Sl'-S\mu
\]

wo $l$ alle Ligaturen der Quadrate in $\omega$ bedeuten (immer so genommen, dass die Quadrate ihnen zur rechten liegen) und wo $l'$ diejenigen Ligaturen bedeutet, die auf dem Umfange der Figur $\omega$ zwischen $0$ und $\frac12m$ liegen, also ausserhalb $\Omega'$.

Alle Ligaturen $l$ bestehen aus
\begin{enumerate}
\item[1)] $l'$
\item[2)] $l''$ die innerhalb $\Omega'$ liegenden Grenzligaturen also $\lambda',\overline{\overline\nu},\overline{\overline\lambda}$.
\item[3)] $l'''$ die im Innern von $\omega$ liegen.
\end{enumerate}
Verstände man unter $l$ indefin. alle Ligaturen, die sich innerhalb $\omega$ oder auf den Grenzen dieser Figur befinden, insofern sie von Punkten $\frac{km}{2M}$ ausgehen, so dass $k$ durch $1+i$ theilbar ist, so wäre der Decident
\[
=\Sigma\alpha\pdot Sl-S\mu
\]

\begin{list}{}{\setlength{\labelwidth}{5em}\setlength{\leftmargin}{5.5em}}
\item[wo $\alpha=1$] für alle Ligaturen im Innern von $\omega$
\item[$\alpha=\Upsilon l+1$] für alle Grenzligaturen ausserhalb $\Omega'$, deren Richtung in der von $0$ nach $\frac12m$ gehenden Grenze liegt
\end{list}

\SourceLeaf{341}

\begin{list}{}{\setlength{\labelwidth}{9em}\setlength{\leftmargin}{9.5em}}
\item[$\alpha=-(\Upsilon\overline{\overline\lambda}+1)=-\Upsilon l$] für alle auf dieser Grenze, die in entgegengesetztem Sinne laufen
\item[$\alpha=\Upsilon l$] für alle Grenzligaturen innerhalb $\Omega'$, deren Richtung auf $0$ zugeht
\item[$\alpha=-\Upsilon l+1$] für alle Grenzligaturen innerhalb $\Omega'$, deren Richtung von $0$ abwärts geht.
\end{list}

\Article{12}
Wir können nun die sämmtlichen vorkommenden $l$ (nach der letzten Manier) zu zweien combiniren, nemlich $l$ mit $\overline{\overline{\frac12m-l}}$, welche wir verbundene Ligaturen nennen wollen; eine einzige ist hiervon ausgenommen, welche isolirt steht oder mit ihrer verbundenen Ligatur identisch ist, nemlich diejenige, welche von
\[
\tfrac12(M-1)\pdot\frac{m}{2M}\quad\text{nach}\quad\tfrac12(M+1)\frac{m}{2M}\quad\text{läuft}
\]
für verbundene Ligaturen ist das $\alpha$ immer einerlei.
\BookRule{.3\textwidth}
